% MT062 / explanation — canonical editable content.
方差衡量预测误差随训练随机性怎样波动；相关系数\(\rho\)衡量两棵树是否倾向一起偏高或一起偏低。树若总以相同方向犯错，平均以后这部分错误仍会留下。


先把题给公式整理成



\[
\operatorname{Var}(F)=\sigma^2\left[\rho+\frac{1-\rho}{n}\right].
\]


\textbf{(a) 增加树数消掉哪一部分？} 固定单树方差与相关性时，\((1-\rho)/n\)随\(n\)增加而减小，但\(\rho\)不变。若\(\rho=0\)，方差按\(1/n\)下降；若\(\rho=1\)，所有树完全一起波动，平均再多也还是\(\sigma^2\)。对固定正\(\rho\)，很多树时的极限为\(\rho\sigma^2\)。


补充数字：\(\sigma^2=1,\rho=0.2\)，5棵树的方差为\(0.2+0.8/5=0.36\)；20棵为\(0.24\)；继续增加时趋于0.2。树数越多，后续新增的收益往往越小。


\textbf{(b) 固定树数怎样改进？} 不同bootstrap样本让树看到不同的训练组合；每个节点随机抽候选特征，避免所有树总在同一特征上切分。这些措施主要降低树之间的误差相关性。也可通过适当控制树深等改变单树波动，但随机化或简化过头会损害单树预测能力，需要一并验证。


选读：设各树误差为\(E_i\)，平均误差为\(\frac1n\sum_iE_i\)。展开方差时有\(n\)个自身方差\(\sigma^2\)，还有\(n(n-1)\)个有序协方差\(\rho\sigma^2\)，总和除以\(n^2\)就得到题式。该公式只描述方差，没有证明偏差、观测噪声或全部测试错误会降到0。


\textbf{出处：} 2023B期中 Q11，原卷第8页（10分）


\textbf{答案依据：} 2023B期中 Q11 随卷／配套参考答案；条件与解释经复核。
